%%%%%%%%%%  scan idx 274  |  folio 261  %%%%%%%%%%
\SGAPrintPage{261}%
of the type \(\underline{O}_{U,X}\), where \(U\) is an object of \(X\). Using the right exactness and commutation with direct sums of the functors \(H^p(E)\otimes\,.\,\) and \(H^p(E\otimes^{\mathrm L}\,.\,)\), we reduce to proving that the arrow of a') is an isomorphism for \(M=\underline{O}_{U,X}\), \(U\in\Ob X\). Now \(H^p(E)\otimes\underline{O}_{U,X}\) (resp. \(H^p(E\otimes^{\mathrm L}\underline{O}_{U,X})\)) identifies canonically (SGA 4 IV 2.13) with \(i_!(H^p(E)|U)\) (resp. \(H^p(i_!(E|U))\)), where \(i_!\) is extension by \(0\) outside \(U\), and the result follows because \(i_!\) is exact.

Suppose \(E\) pseudo-coherent, and let us prove \(\mathrm b)\Longrightarrow\mathrm b')\). It clearly suffices to show that \((*)\) is an isomorphism for \(M\) injective. Since the arrow \((***)\) is an isomorphism for \(M\) injective, the assertion will therefore follow from the following one:

\((5.1.1)\): \uline{the arrow} \((**)\) \uline{is an isomorphism if} \(E\) \uline{is pseudo-coherent and} \(M\) \uline{injective}.

\uline{Proof of} \((5.1.1)\). Denote by \(T\) (resp. \(T'\)) the functor from \(D^-(X)\) to \(D(X)\) defined by \(F\mapsto F\otimes^{\mathrm L}M\) (resp. \(\mathrm R\underline{\Hom}(F^{\vee},M)\)). Let \(n\in\mathbb Z\); if \(H^i(F)=0\) for \(i>n\), then \(H^iT(F)=0\) for \(i>n\), \(H^i(F^{\vee})=0\) for \(i<-n\), and \(H^iT'(F)=0\) for \(i>n\) if \(M\) is injective. Since, on the other hand, the arrow \((**)\) commutes with restriction to an open, and an injective restricted to an open is still injective, we may, after localizing and truncating, suppose \(E\) strictly perfect (I 2.1). By dévissage we then reduce to \(E=\underline{O}_X\), and the result follows.

\sgaStableTarget{sga6.e03.rem.5-2}\textbf{Remarks 5.2.} (i) Condition a) (resp. a'), b), b')) is local in nature: if \((X_\alpha)\) is an open covering of the final object of \(X\), and if, for each \(\alpha\), condition a) (resp\ldots) holds for \(X_\alpha\) and \(E\,|\,X_\alpha\), then a) (resp\ldots) holds.

(ii) If \(x\) is a point of \(X\), denote by \(\mathrm a_x)\) (resp. \(\mathrm a'_x)\), etc.) the condi%
%%%%%%%%%%  scan idx 275  |  folio 262  %%%%%%%%%%
\SGAPrintPage{262}%
tion a) (resp. a'), etc.) relative to \(E_x\) and \(\Mod(\underline{O}_{X,x})\). Then, if \(X\) has enough points, condition a) (resp. a'), b)) is equivalent to the condition ``for every point \(x\) of \(X\), \(\mathrm a_x)\) (resp. \(\mathrm a'_x)\), \(\mathrm b_x))\)'', and is equivalent to the condition ``for every point \(x\) of \(X\), \(\mathrm b'_x)\)'' if \(E\) is pseudo-coherent; moreover, condition a') is equivalent to the analogous condition in which ``for every \(M\in\Ob\Mod(X)\)'' is replaced by ``for every \(\underline{O}_X\)-Module that is the cokernel of a homomorphism \(\underline{O}_{U,X}^p\longrightarrow\underline{O}_{U,X}^q\), where \(U\) is an object of \(X\)''.

(iii) The criterion (EGA III 7.4.2) extends word for word to our situation: if \(E\) has flat components, condition a) simply means that \(\Coker(E^p\longrightarrow E^{p+1})\) is flat, hence locally a direct factor of a free Module of finite type when \(E^p\) and \(E^{p+1}\) are also supposed of finite presentation (I 5.8.3). In particular, we deduce that, if \(E\) is perfect, then \(E\) satisfies condition a) relative to the integer \(p\) if and only if \(E^{\vee}\) satisfies condition a) relative to the integer \(-p-1\).

(iv) Let \(X\) be a \uline{quasi-separated} scheme, \(\underline{O}_X\) the structure sheaf, \(E\in\Ob D^-(X)\) a complex with locally bounded and quasi-coherent cohomology. Denote by \(\mathrm a_{\qcoh})\) (resp. \(\mathrm a'_{\qcoh})\), etc.) the condition analogous to a) (resp. a'), etc.) in which \(\Mod(X)\) is replaced by the category \(\Qcoh(X)\) of quasi-coherent \(\underline{O}_X\)-Modules. Then we have
\[
 \mathrm a)\Longleftrightarrow\mathrm a_{\qcoh})\Longleftrightarrow\mathrm a'_{\qcoh})\Longleftrightarrow\mathrm b_{\qcoh})\Longleftarrow\mathrm b'_{\qcoh}),
\]
and the preceding conditions are equivalent if \(E\) is pseudo-coherent.

Indeed, since \(X\) is quasi-separated, every quasi-coherent sheaf on an affine open of \(X\) is induced by a quasi-coherent sheaf on \(X\) (EGA I 9.4.2 and IV 1.2.7). It follows that condition \(\mathrm a_{\qcoh})\) (resp. \(\mathrm a'_{\qcoh})\), \(\mathrm b_{\qcoh})\), \(\mathrm b'_{\qcoh}))\) holds if and only if it holds over every affine open. To prove the implications above, we may therefore suppose \(X\) affine. Then (II 3.5) \(E\) is isomorphic in \(D(X)\) to a complex of quasi-coherent Modules, and the implications
\[
 \mathrm a_{\qcoh})\Longleftrightarrow\mathrm a'_{\qcoh})\Longleftrightarrow\mathrm b_{\qcoh})\Longleftarrow\mathrm b'_{\qcoh})
\]
%%%%%%%%%%  scan idx 276  |  folio 263  %%%%%%%%%%
\SGAPrintPage{263}%
follow from (5.1) applied to the ``one-point'' topos ringed by \(\Gamma(X,\underline{O}_X)\); the same holds, by (II 2.2.10 a)), for the implication \(\mathrm b_{\qcoh})\Longrightarrow\mathrm b'_{\qcoh})\) when \(E\) is pseudo-coherent. Finally, the equivalence \(\mathrm a)\Longleftrightarrow\mathrm a_{\qcoh})\), or, equivalently, \(\mathrm a')\Longleftrightarrow\mathrm a'_{\qcoh})\) follows from the last remark of (ii) (and the fact that quasi-coherent sheaves ``extend'', since \(X\) is quasi-separated).

\sgaStableTarget{sga6.e03.prop.5-3}\textbf{Proposition 5.3.} (\uline{exchange property}) (EGA III 7.7.5 II). \uline{Let} \(f:X\to Y\) \uline{be a morphism of quasi-compact and quasi-separated schemes}, \(E\in\Ob D^b(X)\) \uline{a complex with quasi-coherent cohomology and finite flat amplitude rel.} \uline{to} \(f\), \uline{and} \(p\in\mathbb Z\). \uline{Consider the following conditions}:
\begin{enumerate}[label=\alph*)]
\item[a)] \uline{The functor} \(M\mapsto R^pf_*(E\otimes_Y^{\mathrm L}M)\) \uline{from} \(\Qcoh(Y)\)\SGAAuthorFootnote{1}{\(\Qcoh(Y)\) = category of quasi-coherent \(\underline{O}_Y\)-Modules.} \uline{to} \(\Qcoh(Y)\)\SGAAuthorFootnote{2}{by (3.7).} \uline{is right exact}.
\item[a')] \uline{The canonical arrow}
\[
 R^pf_*(E)\otimes M\longrightarrow R^pf_*(E\otimes_Y^{\mathrm L}M)
\]
(\uline{the composite of the canonical arrow} \(R^pf_*(E)\otimes M\longrightarrow H^p(Rf_*(E)\otimes^{\mathrm L}M)\) \uline{and the isomorphism} (3.7) \(H^p(Rf_*(E)\otimes^{\mathrm L}M)\xrightarrow{\sim}R^pf_*(E\otimes_Y^{\mathrm L}M)\)) \uline{is an isomorphism for every} \(M\in\Ob\Qcoh(Y)\).
\item[b)] \uline{The functor} \(M\mapsto R^{p+1}f_*(E\otimes_Y^{\mathrm L}M)\) \uline{from} \(\Qcoh(Y)\) \uline{to} \(\Qcoh(Y)\) \uline{is left exact}.
\item[b')] \uline{The canonical arrow}
\[
 R^{p+1}f_*(E\otimes_Y^{\mathrm L}M)\longrightarrow\underline{\mathrm{Hom}}(H^{-p-1}(Rf_*(E)^{\vee}),M)
\]
(\uline{the composite of the isomorphism} (3.7) \(R^{p+1}f_*(E\otimes_Y^{\mathrm L}M)\simeq H^{p+1}(Rf_*(E)\otimes^{\mathrm L}M)\) \uline{and the canonical arrow} (5.1 (\(*\))) \(H^{p+1}(Rf_*(E)\otimes^{\mathrm L}M)\longrightarrow\underline{\mathrm{Hom}}(H^{-p-1}(Rf_*(E)^{\vee}),M)\)) \uline{is an isomorphism for every} \(M\in\Ob\Qcoh(Y)\).
%%%%%%%%%%  scan idx 277  |  folio 264  %%%%%%%%%%
\SGAPrintPage{264}%
\item[c)]
\uline{For every base change} \(Y'\to Y\), \uline{the canonical arrow}
\[
 R^pf_*(E)\otimes_Y\underline{O}_{Y'}\longrightarrow H^p(Rf_*(E)\otimes_Y^{\mathrm L}\underline{O}_{Y'})
\]
\uline{is an isomorphism}.
\end{enumerate}
\uline{Then we have}:
\[
 \mathrm a)\Longleftrightarrow\mathrm a')\Longleftrightarrow\mathrm b)\Longleftrightarrow\mathrm c)\Longleftarrow\mathrm b').
\]
\uline{Moreover}, \uline{if} \(f\) \uline{is proper}, \(E\) \uline{pseudo-coherent rel.} \uline{to} \(f\) (\uline{hence perfect rel.} \uline{to} \(f\)), \uline{and subject to the validity of} (2.1) \uline{for} \(f\) (\uline{hence in particular if} \(f\) \uline{is projective or if} \(Y\) \uline{is noetherian}), \uline{the preceding conditions are equivalent}.

\uline{Proof}. The implications \(\mathrm a)\Longleftrightarrow\mathrm a')\Longleftrightarrow\mathrm b)\Longleftarrow\mathrm b')\) follow at once from (3.7), (5.1) and (5.2 (iv)), as does the implication \(\mathrm b')\Longleftarrow\mathrm b)\) (for \(f\) proper and \(E\) pseudo-coherent rel. to \(f\)) by (2.1). It remains to prove that \(\mathrm a)\Longleftrightarrow\mathrm c)\). This will follow from the first assertion of the following lemma:

\sgaStableTarget{sga6.e03.lem.5-3-1}\textbf{Lemma 5.3.1.} \uline{Let} \(X\) \uline{be a quasi-separated scheme}, \uline{and let} \(E\in\Ob D^{-}(X)\) \uline{be a complex with locally bounded and quasi-coherent cohomology}. \uline{Then condition} (5.1 a)) \uline{is equivalent to the condition}
\begin{enumerate}[label=\alph*)]
\item[c)] \uline{For every scheme} \(X'\) \uline{over} \(X\), \uline{the canonical arrow}
\[
 H^p(E)\otimes_X\underline{O}_{X'}\longrightarrow H^p(E\otimes_X^{\mathrm L}\underline{O}_{X'})
\]
\uline{is an isomorphism}.
\end{enumerate}
\uline{Moreover}, \uline{if} \(E\) \uline{is pseudo-coherent}, \uline{condition} (5.1 a)) \uline{implies that}, \uline{for every morphism} \(X'\to X\), \uline{the canonical arrow}
\sgaStableTarget{sga6.e03.eq.5-3-1-1}\begin{equation}
 H^{-p-1}(E^{\vee})\otimes_X\underline{O}_{X'}\longrightarrow H^{-p-1}((E\otimes_X^{\mathrm L}\underline{O}_{X'})^{\vee}),\tag{5.3.1.1}
\end{equation}
%%%%%%%%%%  scan idx 278  |  folio 265  %%%%%%%%%%
\SGAPrintPage{265}%
\uline{the composite of} \(H^{-p-1}(E^{\vee})\otimes_X\underline{O}_{X'}\longrightarrow H^{-p-1}(E^{\vee}\otimes_X^{\mathrm L}\underline{O}_{X'})\) \uline{and the canonical arrow} (I 7.1.2) \(H^{-p-1}(E^{\vee}\otimes_X^{\mathrm L}\underline{O}_{X'})\longrightarrow H^{-p-1}((E\otimes_X^{\mathrm L}\underline{O}_{X'})^{\vee})\), \uline{is an isomorphism}.

\uline{Proof}. \(\mathrm a)\Longrightarrow\mathrm c)\). The question being local on \(X'\), we may suppose \(X\) (resp. \(X'\)) affine with ring \(A\) (resp. \(A'\)). By (II 3.5), we may also suppose that \(E=\widetilde F\), where \(F\) is a complex of \(A\)-modules. By hypothesis (condition a')), the arrow
\[
 H^p(F)\otimes_A A'\longrightarrow H^p(F\otimes_A^{\mathrm L}A')
\]
is an isomorphism, and the result follows since \(H^p(E)\otimes_X\underline{O}_{X'}\) (resp. \(H^p(E\otimes_X^{\mathrm L}\underline{O}_{X'})\)) identifies canonically with the sheaf on \(X'\) defined by the \(A'\)-module \(H^p(F)\otimes_A A'\) (resp. \(H^p(F\otimes_A^{\mathrm L}A')\)).

\(\mathrm c)\Longrightarrow\mathrm a)\). In view of (5.2 (iv)), it suffices to prove \(\mathrm c)\Longrightarrow\mathrm a'_{\qcoh})\). The question being local (since \(X\) is quasi-separated), we may suppose \(X\) affine with ring \(A\), and \(E=\widetilde F\), where \(F\) is a complex of \(A\)-modules. Let \(M=\widetilde N\) be a quasi-coherent \(\underline{O}_X\)-Module. We must show that the arrow
\[
 H^p(F)\otimes_A N\longrightarrow H^p(F\otimes_A^{\mathrm L}N)
\]
is an isomorphism. But, putting \(A'=D_A(N)\) (EGA \(\mathrm O_{\mathrm{IV}}\) 18.2.3), \(X'=\Spec(A')\), this follows from c) applied to the canonical morphism \(X'\to X\).

Let us prove the last assertion. The question being local on \(X'\), we may, by (II 2.2.10 b)), suppose \(E\) strictly pseudo-coherent. Moreover, since the assertion involves only a part of \(E\) of bounded length, we may, after truncating, suppose \(E\) strictly perfect. Then \((E\otimes_X^{\mathrm L}\underline{O}_{X'})^{\vee}\) identifies canonically (I 7.1.2) with \(E^{\vee}\otimes_X^{\mathrm L}\underline{O}_{X'}\), and the arrow (5.3.1.1) may be rewritten as the canonical arrow
\sgaStableTarget{sga6.e03.eq.unnamed.occ2}\begin{equation}
 H^{-p-1}(E^{\vee})\otimes_X\underline{O}_{X'}\longrightarrow H^{-p-1}(E^{\vee}\otimes_X^{\mathrm L}\underline{O}_{X'}).\tag{$*$}
\end{equation}
